\chapter*{INTRODUCTION: BASIC CONCEPTS}
The evaluation of many physical quantities, such as expectation values of energy, electric and magnetic multipole moments, transition probabilities, etc., is considerably simplified by making use of the transformation properties of these quantities under coordinate rotation and inversion in three-dimensional space.

The transformation properties of physical quantities with respect to rotations reveal themselves either through rotations of the given physical system relative to some fixed reference frame or through rotations of the coordinate axes relative to the physical system.

Inversion characterizes the behavior of a physical quantity under transformation from a right-handed coordinate system to a left-handed one, or vice versa.

Many physical quantities, by their nature, are invariants under coordinate rotations. In particular, the properties of any closed physical system should be independent of rotations, as follows from the isotropy of space. As a consequence of this fundamental property of space, the total angular momentum of such a system is an integral of motion.

A similar situation occurs with regard to coordinate inversion. Excluding phenomena connected with the weak interaction, a wealth of atomic, molecular and nuclear processes look alike in right-handed and left-handed coordinate systems. As a result of this "mirror" symmetry, quantum states of atoms, molecules, nuclei and elementary particles may be characterized by definite parity.

Strictly speaking, a quantum mechanical wave function $\Psi(r)$ of any closed physical system may be characterized by four quantum numbers ( $\varepsilon, \pi, j, m$ ) which are the eigenvalues of four commuting operators: the Hamiltonian $\widehat{H}$, the parity operator $\widehat{P}_{\tau}$, the operator $\widehat{\vect{J}}^{2}$ of the square of the angular momentum and the operator $\widehat{J}_{z}$ of the projection of this momentum onto a quantization axis. Thus $\Psi(r) \equiv \Psi_{\varepsilon \pi \alpha j m}(r)$ obeys the equations

\[
\begin{gathered}
\widehat{H} \Psi_{\varepsilon \pi \alpha j m}(r)=\varepsilon \Psi_{\varepsilon \pi \alpha j m}(r), \\
\widehat{P}_{r} \Psi_{\varepsilon \pi \alpha j m}(r)=\pi \Psi_{e \pi \alpha j m}(r), \\
\widehat{\vect{J}}^{2} \Psi_{e \pi \alpha j m}(r)=j(j+1) \Psi_{\varepsilon \pi \alpha j m}(r), \\
\widehat{J}_{z} \Psi_{\varepsilon \pi \alpha j m}(r)=m \Psi_{\varepsilon \pi \alpha j m}(r),
\end{gathered}
\]

where $\alpha$ denotes all other quantum numbers (if available) and $r$ represents a variety of arguments.\\
$\varepsilon, \pi, j$ and $m$ are the integrals of motion not only for a closed system, but also for any system exposed to the action of an external spherically-symmetric field. Moreover, the ( $\varepsilon \pi j m$ )-representation appears to be very convenient for practical calculation even if the external field is not spherically-symmetric.

For a given $j$, there exist $2 j+1$ wave functions which correspond to different $m$. These functions describe the quantum states of the system which differ only by the orientation of the angular momentum in space.

Under coordinate rotations these functions undergo linear mutual transformations which do not involve the functions of other quantum numbers ( $\varepsilon, \pi, \alpha, j$ ) and may be written as

\[
\Psi_{c \pi \alpha j m^{\prime}}\left(r^{\prime}\right)=\sum_{m} \Psi_{c \pi \alpha j m}(r) D_{m m^{\prime}}^{j}(\alpha, \beta, \gamma)
\]

where the transformations coefficients $D_{m m^{\prime}}^{j}(\alpha, \beta, \gamma)$ are called Wigner $D$-functions. They are the elements of the finite rotation matrix in the $j$-representation and depend on the Euler angles $\alpha, \beta, \gamma$ which determine rotation.

Let a quantum-mechanical system which possesses some fixed angular momentum $j$ and its projection $m$ consist of two subsystems, with certain angular momenta $j_{1}$ and $j_{2}$, respectively. In this case the wave function $\Psi_{j_{1} j_{2} j m}\left(r_{1}, r_{2}\right)$ may be constructed from the wave functions of subsystems according to the relation

\[
\Psi_{j_{1} j_{2} j m}\left(r_{1}, r_{2}\right)=\sum_{m_{1} m_{2}} \clebsch{j_{1}}{m_{1}}{j_{2}}{m_{2}}{j}{m} \Psi_{j_{1} m_{1}}\left(r_{1}\right) \Psi_{j_{2} m_{2}}\left(r_{2}\right)
\]

The quantities $\clebsch{j_{1}}{m_{1}}{j_{2}}{m_{2}}{j}{m}$ are called the Clebsch-Gordan coefficients. They are very important in quantum mechanics because they allow one to construct wave functions for various complex systems (nuclei, atoms, molecules, etc.). The addition of many angular momenta into some resultant total angular momentum may be performed in different ways, or, in other words, in accordance with different coupling schemes of the momenta in question. Each scheme may be associated with a certain representation of these angular momenta. The unitary transformations which relate various representations and describe the recoupling of angular momenta are realized by matrices expressed in terms of the $6 j$ - $9 j$-symbols and others $3 n j$-symbols of higher order.

One of the principal concepts of the quantum theory of angular momentum is that of an irreducible tensor. By definition, an irreducible tensor of rank $\lambda$ has $2 \lambda+1$ components which transform under coordinate rotation as

\[
\mathfrak{M}_{\lambda \mu}=\sum_{\mu^{\prime}} \mathfrak{M}_{\lambda \mu^{\prime}} D_{\mu^{\prime} \mu}^{\lambda}(\alpha, \beta, \gamma)
\]

In particular, the wave functions $\Psi_{j m}$ with fixed $j$ but different $m$ constitute an irreducible tensor of rank $j$. Moreover, irreducible tensors may be constructed from various physical quantities. For instance, energy is a tensor of zero rank (scalar), spin and magnetic moments are tensors of first rank (vectors), the quadrupole moment is a tensor of second rank, etc. Generally, any physical quantity or operator which corresponds to these quantities may be represented as a linear combination of irreducible tensors.

The introduction of irreducible tensors into the analysis of physical quantities or corresponding operators substantially simplifies the evaluation of matrix elements of irreducible tensor operators,

\[
\braOket{\varepsilon^{\prime} \pi^{\prime} \alpha^{\prime} j^{\prime} m^{\prime}}{\widehat{\mathfrak{M}}_{\lambda \mu}}{\varepsilon \pi \alpha j m}=\int \Psi_{c^{\prime} \pi^{\prime} \alpha^{\prime} j^{\prime} m^{\prime}}^{*}(r) \widehat{\mathfrak{M}}_{\lambda \mu} \Psi_{c \pi \alpha j m}(r) d r
\]

where $\int \ldots d r$ denotes integration over continuous variables and summation over discrete ones. These matrix elements determine expectation values of physical quantities in definite quantum states, the probabilities of quantum transitions between various states, etc.

According to the Wigner-Eckart theorem, all the dependence of matrix elements on the orientation of coordinate axes, i.e., on quantum numbers $m, m^{\prime}, \mu$ which determine the projections of angular momenta, is entirely contained in the Clebsch-Gordan coefficient

\[
\braOket{\varepsilon^{\prime} \pi^{\prime} \alpha^{\prime} j^{\prime} m^{\prime}}{\widehat{\mathfrak{M}}_{\lambda \mu}}{\varepsilon \pi \alpha j m}=\clebsch{j}{m}{\lambda}{\mu}{j^{\prime}}{m^{\prime}}\braOketred{\varepsilon^{\prime} \pi^{\prime} \alpha^{\prime} j^{\prime}}{\widehat{\mathfrak{M}}_{\lambda}}{\varepsilon \pi \alpha j} .
\]

In this case $\braOketred{\varepsilon^{\prime} \pi^{\prime} \alpha^{\prime} j^{\prime}}{\widehat{\mathfrak{M}}_{\lambda}}{\varepsilon \pi \alpha j}$ is an invariant factor called a reduced matrix element. The Wigner-Eckart theorem allows one to reduce evaluation of transition probabilities, angular distributions, polarizations, etc., to the calculation of standard sums of the vector addition coefficients and recoupling coefficients.

All the facts mentioned above reveal that the quantum theory of angular momentum provides the formalism which is universal and extremely convenient for various practical calculations.


